% from ega4-2-fr.tex
\oldpage[IV-2]{52}

\section{Changement du corps de base dans les préschémas algébriques}
\label{IV.2.4-fr}

\subsection{Dimension des préschémas algébriques}
\label{subsection:IV.2.4.1-fr}

Nous développerons au \S~5 la théorie générale de la dimension des préschémas~;
mais la théorie de la dimension des préschémas \emph{algébriques} peut se
développer de façon plus élémentaire, et comme elle présente en outre de
nombreux caractères spéciaux, nous en donnerons ici un exposé rapide,
indépendant de la théorie générale.

Si $K$ est un corps et $L$ une extension de $K$, on notera
$\operatorname{deg.tr.}_K L$ le \emph{degré de transcendance} de $L$ sur $K$.

\begin{definition}[4.1.1]
\label{IV.2.4.1.1-fr}
Soient $k$ un corps, $X$ un préschéma localement de type fini sur $k$. On
appelle \emph{dimension} de $X$ le nombre
\begin{equation}
\tag{4.1.1.1}
\label{IV.2.4.1.1.1-fr}
\dim X=\sup_x\operatorname{deg.tr.}_k k(x)
\end{equation}
où $x$ parcourt l'ensemble des points maximaux de $X$.
\end{definition}

Nous verrons au \S~5 (5.2.2) que la définition (4.1.1) ne dépend qu'en
apparence du corps de base $k$ sur lequel $X$ est localement de type fini, et
que le nombre $\dim(X)$ ainsi défini coïncide avec la dimension de
l'\emph{espace sous-jacent} $X$ (définie dans $(0,14.1.2)$). On a évidemment
$\dim X=\dim X_{\mathrm{red}}$.

On notera que chaque $k(x)$ est une extension \emph{de type fini} de $k$
(I, 6.3.3) donc a un degré de transcendance fini sur $k$. Si $X$ est de type
fini sur $k$ et non vide, il est noethérien (I, 6.3.7), donc n'a qu'un nombre
fini de composantes irréductibles, et par suite $\dim(X)$ est \emph{finie} et
$\geqslant0$~; la définition (4.1.1) donne
\[
\dim(\varnothing)=-\infty.
\]

Il est clair que, si l'on désigne par $(X_\alpha)$ la famille des
sous-préschémas fermés réduits de $X$ ayant pour espaces sous-jacents les
composantes irréductibles de $X$ (I, 5.2.1), on a
\begin{equation}
\tag{4.1.1.2}
\label{IV.2.4.1.1.2-fr}
\dim(X)=\sup_\alpha\dim(X_\alpha).
\end{equation}
Cela ramène donc le calcul de la dimension au cas des préschémas \emph{intègres}
(localement de type fini sur $k$). Enfin, on a évidemment
\begin{equation}
\tag{4.1.1.3}
\label{IV.2.4.1.1.3-fr}
\dim(X)=\dim(U)
\end{equation}
pour tout sous-préschéma induit sur un ouvert \emph{partout dense} $U$ de
$X$~; cela ramène donc finalement la notion de dimension au cas d'un schéma
\emph{affine de type fini sur $k$}.

\begin{theorem}[4.1.2]
\label{IV.2.4.1.2-fr}
Soient $X$, $Y$ deux préschémas localement de type fini sur un corps $k$,
$f:X\to Y$ un $k$-morphisme.

